\section*{G.1 Overview}

This appendix summarises the ethical and practical risks considered during the design and conduct of this dissertation, and the mitigations applied.

\section*{G.2 Risk: Use of Sensitive Mental Health Content}

\textbf{Risk.} The dissertation's subject matter involves mental health crisis scenarios, including scenarios consistent with an implied risk of self-harm, which could be distressing to encounter during data collection, human coding, or dissemination.

\textbf{Mitigation.} No real user disclosures were used at any stage: all scenarios are synthetic constructions written for this dissertation (Appendix A), and all response data is model-generated text produced in response to these synthetic prompts. No real patients, clinical cases, or crisis-service call transcripts were accessed or used. Human coders labelled model-generated text under a structured coding protocol (Appendix B.3--B.4) rather than open-ended exposure to unstructured crisis content.

\section*{G.3 Risk: No Real-World Intervention or Clinical Contact}

\textbf{Risk.} Because the scenarios describe crisis-adjacent situations, there is a theoretical risk of conflating this research activity with the provision of real mental health support, or of a reader mistaking the dissertation's synthetic scenarios for real cases requiring intervention.

\textbf{Mitigation.} This dissertation is an audit of AI system behaviour, not a mental health intervention study; at no point does it involve, solicit, or respond to real disclosures of distress. This is stated explicitly in the Methodology (Section 3.6) and Electronic Supplementary Material (Appendix E.4) sections.

\section*{G.4 Risk: Reproduction of Harmful or Fabricated Content}

\textbf{Risk.} The crisis-contact fabrication analysis (Section 4.7) necessarily involves identifying and, in Appendix D.4, reproducing examples of fabricated crisis-helpline numbers generated by the models under study. There is a risk that reproducing such numbers, even as illustrative examples of a fabrication pattern, could inadvertently propagate a non-functional number if excerpted out of context.

\textbf{Mitigation.} Examples reproduced in Appendix D.4 are presented explicitly and exclusively as illustrations of a fabrication \emph{pattern} (e.g. the North American fictional ``555'' exchange, sequential digit runs), with the fabricated status of each number stated directly alongside it, and are not presented as, or capable of being mistaken for, genuine crisis-line guidance. No real, currently operative crisis-line numbers are altered or obscured in this reporting.

\section*{G.5 Risk: Cultural and Religious Sensitivity}

\textbf{Risk.} H3 concerns the framing of religious and spiritual support by geographic region, a topic requiring careful handling to avoid the dissertation itself making essentialising claims about the religiosity of any region or population.

\textbf{Mitigation.} The dissertation's claims are restricted explicitly to the behaviour of the language models under study -- what they recommend, and to whom -- and do not make any claim about the actual religious composition, preferences, or needs of populations in any region. The Discussion (Section 5.3) explicitly notes the normative ambiguity of this finding and declines to characterise the observed pattern as straightforwardly harmful without further, qualitative user research.

\section*{G.6 Risk: Data Privacy}

\textbf{Risk.} None identified. No personal data of any kind -- from human coders, from real users, or from any other individual -- was collected, stored, or processed at any stage of this research. Human coders labelled anonymised, model-generated response text only; no coder-identifying information is retained in any dataset produced.

\section*{G.7 Risk: Researcher Wellbeing}

\textbf{Risk.} Repeated exposure to crisis-related content during data collection, coding, and analysis could affect researcher wellbeing over the course of the project.

\textbf{Mitigation.} All content reviewed was synthetic and model-generated rather than real crisis disclosures, substantially reducing the emotional burden relative to work involving real clinical or crisis-service data. The coding protocol (Appendix B.3--B.4) structured the task as a defined labelling exercise against explicit written criteria, rather than open-ended review, and the volume of content per session was limited by the stratified sampling design (112 and 100 responses respectively, rather than the full 1{,}120-response dataset).
